\section{Claude Fable 5.2とMythos 5.2を発表}
AnthropicはClaude Fable 5.2とMythos 5.2を発表した。Grokの収集によれば、コーディング、推論、画像理解の改善に加え、出力速度を最大30\%高め、入力100万トークン当たり8ドル、出力40ドルへ価格を引き下げた。100万トークンの文脈長と最大12万8000トークンの出力にも対応するという。Fableは一般提供、Mythosは検証プログラム経由とされる。形式化数学、ロボット制御、物理研究の成果も紹介されたが、評価の詳細は公式発表に基づく主張として読む必要がある。
\dailyxsource{Claude公式投稿}{https://x.com/claudeai/status/2119867114450268322}

\section{Metaが高効率モデルMuse Emberを公開}
Meta Superintelligence Labsは、ツール呼び出しやコード処理、情報抽出など、エージェントの反復的な作業向けにMuse Emberを発表した。テキスト・画像・音声・動画を扱い、文脈長は100万トークンとされる。入力100万トークン当たり0.10ドル、出力0.50ドルという価格を提示。Artificial Analysisの指数で36点、Haiku 5.5と同程度の水準を約3分の1の費用で実現したとの説明もある。Google Cloudでの公開プレビューも告知された。実際のワークフローでの費用対効果が焦点となる。
\dailyxsource{Meta公式投稿}{https://x.com/AIatMeta/status/2119853591703998787}

\section{Seedream 5.5、連続画像編集の安定性を訴求}
ByteDance Seedは画像生成・編集モデルSeedream 5.5を発表した。最大100回の連続編集でも編集対象外の領域をピクセルレベルで維持できること、透明背景を扱えるネイティブRGBA出力、光や素材の表現、小さな文字の描画改善を特徴としている。即夢、豆包、火山方舟などで展開されるという。2K画像編集の価格や他社モデルを上回るとの比較も投稿されているが、これらの数値や優位性は別途検証が必要だ。画像制作における再編集の再現性が注目される。
\dailyxsource{ByteDance Seed公式投稿}{https://x.com/ByteDanceSeed_/status/2119720475068920181}

\section{Mistral Studioのエージェント構築機能}
Mistral AIはStudio上で、AppsとManaged Agentsを公開プレビューとして、Workflowsを一般提供として利用できることを改めて案内した。Managed Agentsでは長時間の自律実行や障害からの復旧を支援するとされ、企業向けのアプリ・エージェント構築を一体化する方向がうかがえる。機能の一部は観測窓より前に発表されており、今回は当日の公式投稿による再周知を記録する。今後は利用制限や料金、運用事例が重要になる。
\dailyxsource{Mistral公式投稿}{https://x.com/MistralAI/status/2119813353363370134}

\section{次期フロンティアモデルの公開予測}
X上では、GPT-6.2 Astra／Sol、Fable 5.5、Gemini 4 Carbon、Qwen 4、DeepSeek V4.1 Proなどを「今週公開予定」とする投稿が複数流通した。Grokはこれらを未確認と判定しており、公式発表や一次資料に裏付けられた日程ではない。内部テストやリークを示唆する投稿もあるが、確定した製品計画として扱うことはできない。正式な公開情報とコミュニティの期待を区別して追う必要がある。
\dailyxsource{観測窓内の予測投稿}{https://x.com/whooisaakash/status/2119844488806060424}

\section{OpenAIのdotで複数PC接続に言及}
OpenAIの公式アカウントは、複数のローカルコンピュータを同一のdotに接続し、タスクに応じて利用するコンピュータを選べるようになったという開発者向け投稿を共有した。複数端末をまたぐエージェント実行環境として関心を集めるが、Grok原文では対応OS、権限管理、提供条件などの詳細は確認されていない。今回の観測では機能更新の告知を記録するにとどめる。
\dailyxsource{OpenAI公式アカウントの共有}{https://x.com/OpenAI/status/2119889726043992302}

\section{Fable 5.2の外部サービス統合}
Higgsfield AIはClaude Fable 5.2にMCP経由で対応し、3日間の無料・無制限アクセスを提供すると投稿した。Blenderを使った3Dモデリングのデモも紹介され、新モデルのコンピュータ操作やコーディング能力を実務フローに組み込む動きとして注目される。ただし、提供条件や制限の独立確認は本号では行っていない。モデル公開直後のプラットフォーム側の対応速度を示す事例である。
\dailyxsource{Higgsfield公式投稿}{https://x.com/higgsfield_ai/status/2119884452054937656}

\section{Voxtral Transcribe 3のリアルタイム実演}
Mistral AIはAI Engineer Parisで、Voxtral Transcribe 3を用いたリアルタイム文字起こしを実演した。文脈に応じたバイアス設定により認識精度を高める機能が紹介されている。会議や専門領域の音声認識において、事前に与える語彙や文脈の扱いが実用性を左右する。Grok原文では詳細なモデル仕様や定量評価は示されておらず、今回は公式投稿による実演の観測として扱う。
\dailyxsource{Mistral公式投稿}{https://x.com/MistralAI/status/2119822323838832777}